% What FSDS does for continuous inference of multimodal LMs.
% Sits after Table~\ref{tab:msrvtt-mm-shift-contrib}.

\paragraph{What FSDS justifies in continuous multimodal inference.}
Continuous inference for a video--audio--language model is a streaming decision problem: at each window $t$ the system must choose which tokens to re-encode, which KV entries to invalidate, which towers to adapt, and which evidence to write into the belief state.
FSDS does not attribute next-token gradients of a frozen LLM; it attributes nonstationarity in the \emph{input measure} $P(X\mid W)$ on the same $768$/$512$/$768$ blocks that become those tokens (Table~\ref{tab:msrvtt-mm-shift-contrib}).
That is the right object for scheduling continuous inference.
If a modality is stationary across $W$, re-encoding it every window is wasted compute and a shortcut the language head can latch onto.
If it is nonstationary, the KV cache and the reasoner's working memory are stale unless that stream is refreshed.

Concretely, the estimated shares $\pi_{\mathrm{v}}\approx 0.77$--$0.90$, $\pi_{\mathrm{a}}\approx 0.10$--$0.19$, $\pi_{\mathrm{t}}\approx 0$ imply four inference-time controls.
\emph{(i)~Encode-once versus encode-every-window.}
Clip-level text is window-invariant on this extract (mean $|d|=0$), so caption tokens should be encoded once per clip and reused.
Visual tokens, and secondarily audio tokens, must be re-encoded on each window: that is where $P(X\mid W)$ actually moves.
\emph{(ii)~KV invalidation and memory writes.}
A streaming VLM should drop or compress visual (and audio) KV when region-wise $|d|$ or RF-Domain AUC on $W$ crosses a video-clustered threshold, and should \emph{not} refresh caption KV on the same trigger.
A global pooled $d$ is the wrong trigger: signed region effects cancel across videos, so a session-level (video-clustered) monitor matches the sampling unit of the Friedman/Wilcoxon tests.
\emph{(iii)~Test-time adaptation mask.}
If continuous inference includes a test-time LoRA or gated residual, the update support should be the vision tower and fusion, then audio, not the language tower against $W$.
Step sizes follow the same adapt-for-distribution-shift map $\eta_m=\eta_0\pi_m^{\mathrm{RF}}$: more shift mass, larger step; shrinking mass, smaller step.
A rising text VIMP mid-stream is a leakage diagnostic (the model is fitting the static caption), not a reason to spend test-time compute on text.
\emph{(iv)~Belief-state and chain-of-thought evidence.}
For streaming VQA or long-horizon reasoning (``what is happening \emph{now}''), FSDS ranks the evidence that is allowed to revise the belief: visual first, audio second, clip-level subtitle never---unless a window-aligned ASR stream is substituted and re-tested.
Video~13 being audio-dominant further implies that fusion gates must be session-adaptive rather than a frozen $(\pi_{\mathrm{v}},\pi_{\mathrm{a}},\pi_{\mathrm{t}})$.

The justified claim is therefore operational, not metaphysical: FSDS supplies the \emph{nonstationarity budget} of the multimodal stream.
Continuous inference should spend that budget on vision and audio, keep clip-level text as a control variate, and trigger refreshes at video granularity.
What FSDS does \emph{not} yet justify is a claim about which LLM layers implement that budget; that is the next-round streaming result the shares are designed to prioritize.
